% Präambel
\documentclass[12pt,a4paper,oneside, 
liststotoc, 					% Tabellen- und Abbildungsverzeichnis ins Inhaltsverzeichnis
bibtotoc,						% Literaturverzeichnis ins Inhaltsverzeichnis aufnehmen
titlepage, 						% Titlepage-Umgebung statt \maketitle
headsepline, 					% horizontale Linie unter Kolumnentitel
%abstracton,					% Überschrift beim Abstract einschalten, Abstract muss dazu in {abstract}-Umgebung stehen
%DIV11,							% auskommentieren, um den Seitenspiegel zu vergrößern
BCOR6mm,						% Bindekorrektur, die den Seitenspiegel um 6mm nach rechts verschiebt,
]{scrreprt}			
\usepackage{ucs} 				% Dokument in utf8-Codierung schreiben und speichern
\usepackage[utf8x]{inputenc} 	% ermöglicht die direkte Eingabe von Umlauten
\usepackage[ngerman]{babel} 	% deutsche Trennungsregeln und Übersetzung der festcodierten Überschriften
\usepackage[T1]{fontenc} 		% Ausgabe aller zeichen in einer T1-Codierung (wichtig für die Ausgabe von Umlauten!)
\usepackage{graphicx}  			% Einbinden von Grafiken erlauben
\usepackage{amsmath}
\usepackage{amsfonts}
\usepackage{amssymb}
\usepackage{mathpazo} 			% Einstellung der verwendeten Schriftarten
\usepackage{textcomp} 			% zum Einsatz von Eurozeichen u. a. Symbolen
\usepackage{listings}			% Datstellung von Quellcode mit den Umgebungen {lstlisting}, \lstinline und \lstinputlisting
\usepackage{xcolor} 			% einfache Verwendung von Farben in nahezu allen Farbmodellen
\usepackage[intoc]{nomencl}  	% zur Erstellung des Abkürzungsberzeichnisses
\usepackage{fancyhdr}			% Zusatzpaket zur Gestaltung von Fuß und Kopfzeilen
\usepackage{titleref}			% Zum referenzieren mit Überschrift
\usepackage{varwidth}			%alighn asci figure
\usepackage{chngcntr}			% Tabellen- und Abbildungsverzeichnis als fortlaufende Nummer
\usepackage{hyperref}
\usepackage[justification=centering]{caption} % Bild unterschrift mittig ausrichten
\usepackage{nameref}
\counterwithout{figure}{chapter}
\counterwithout{table}{chapter}
\usepackage{url}
\usepackage{ulem}


%For flow chart
\usepackage{tikz}
\usetikzlibrary{shapes,arrows}
% Define block styles
\tikzstyle{decision} = [diamond, draw, fill=blue!20, 
    text width=4.5em, text badly centered, node distance=3cm, inner sep=0pt]
\tikzstyle{block} = [rectangle, draw, fill=blue!20, 
    text width=5em, text centered, rounded corners, minimum height=4em, node distance=3cm]
\tikzstyle{line} = [draw, -latex']
\tikzstyle{cloud} = [draw, ellipse,fill=red!20, node distance=2cm,
    minimum height=2em]


%\usepackage[none]{hyphenat} %deaktiviert Silbentrennung
%\usepackage{showframe}% zum Anzeigen des Seitenlayouts
%Abstand vor Chapter
\renewcommand*\chapterheadstartvskip{\vspace*{-\topskip}}
\renewcommand*\chapterheadendvskip{%
  \vspace*{1\baselineskip plus .1\baselineskip minus .167\baselineskip}}


\setlength\abovedisplayshortskip{0pt}
\setlength\belowdisplayshortskip{0pt}
\setlength\abovedisplayskip{20pt}
\setlength\belowdisplayskip{20pt}

% -----------------------------------------------------------------------------------------------------------------
% Zum Aktualisieren des Abkürzungsverzeichnisses bitte auf der Kommandozeile folgenden Befehl aufrufen :
%  makeindex Bachelorarbeit.nlo -s nomencl.ist -o Bachelorarbeit.nls
% -----------------------------------------------------------------------------------------------------------------

% Hier die persönlichen Daten eingeben:

\newcommand{\titel}{Titel}
\newcommand{\untertitel}{Projekt untertitel}
\newcommand{\arbeit}{Projektarbeit}
\newcommand{\studiengang}{Informatik}
\newcommand{\autor}{Maximilian Bauknracht\\ Nina Zaske\\ Philippe Käufer}
\newcommand{\kurs}{TINF15K}
\newcommand{\abgabe}{\today}
\newcommand{\betreuerdhbw}{Maximilian Scholl}

\newcommand{\jahr}{2017 - 2018}			% für Angabe im Copyright-Vermerk der Titelseite

% Abkürzungen
\newcommand{\ua}{\mbox{u.\,a.\ }}
\newcommand{\zB}{\mbox{z.B.\ }}
\newcommand{\bs}{$\backslash$}
\newcommand{\Csharp}{C\#}

\renewcommand{\nomname}{Abkürzungsverzeichnis}
\makenomenclature 

% -------------------------------------------------------------------------------------------
% Definition der Kopf- und Fußzeilen
\lhead{}								% Kopf links
\chead{}								% Kopf mitte
\rhead{\sffamily{\titel}}				% Kopf rechts
\lfoot{}								% Fuß links
\cfoot{\sffamily{\thepage}}				% Fuß mitte
\rfoot{\sffamily{\autor}}				% Fuß rechts
\renewcommand{\headrulewidth}{0.4pt}	% Liniendicke Kopf
\renewcommand{\footrulewidth}{0.4pt}	% Liniendicke Fuß

\makenomenclature						% Abkürzungsverzeichnis erstellen

% alle Abkürzungen, die in der Arbeit verwendet werden

\nomenclature{DHBW}{Duale Hochschule Baden-Württemberg}
				% Datei mit Abkürzungen laden

%%makeindex t2000.nlo -s nomencl.ist -o t2000.nls %% Abkürzungsverzeichnis erstellen

% -------------------------------------------------------------------------------------------
%                     Beginn des Dokumenteninhalts
% -------------------------------------------------------------------------------------------
\begin{document}
\setcounter{secnumdepth}{2}					% Nummerierungstiefe fürs Inhaltsverzeichnis
\setcounter{tocdepth}{2}
\sffamily									% für die Titelei serifenlose Schrift verwenden

% ------------------------------ Titelei -----------------------------------------------------

\thispagestyle{plain}
\begin{titlepage}
\enlargethispage{4.0cm}
%\sffamily 								% Serifenlose Grundschrift für die Titelseite einstellen
\begin{figure}[!htb]
\hfill
\minipage{0.4\textwidth}%
  \includegraphics[width=\linewidth]{Bilder/logo_dhbw.jpg}
\endminipage
\end{figure}

\begin{center}

\huge{\textsc{\textbf{\titel}}}\\[1.5ex]
\Large{\textbf{\untertitel}}\\[5ex]
\LARGE{\textbf{\arbeit}}\\[2ex]
\Large{Studiengang \studiengang}\\[1ex]
\normalsize{Bachelor of Science}\\[1ex]
\normalsize{Duale Hochschule Baden-Württemberg Stuttgart}\\[3ex]

von\\[1ex] \LARGE{\autor} \\[5ex]


\end{center}

\begin{flushleft}

\begin{tabular}{ll}
Abgabedatum:					& \quad \abgabe \\
%Bearbeitungszeitraum:			& \quad 12 Wochen   \\ 
Kurs: 			& \quad \kurs \\ 
Betreuer der Dualen Hochschule: & \quad \betreuerdhbw \\ [5ex]

\end{tabular} 



\small
Copyrightvermerk:\\

Dieses Werk einschließlich seiner Teile ist \textbf{urheberrechtlich geschützt}. Jede Verwertung außerhalb der engen Grenzen des Urheberrechtgesetzes ist ohne Zustimmung des Autors unzulässig und strafbar. Das gilt insbesondere für Vervielfältigungen, Übersetzungen, Mikroverfilmungen sowie die Einspeicherung und Verarbeitung in elektronischen Systemen.
\end{flushleft}
\begin{flushright}
\copyright{} \jahr
\end{flushright}
\end{titlepage}
\rmfamily 				% erzeugt die Titelseite
\pagenumbering{Roman}						% große, römische Seitenzahlen für Titelei
\addchap{Eidesstattliche Erklärung}
Ich versichere hiermit, dass ich meine Praxisarbeit mit dem Thema
\begin{quote}
\textit{\titel} -\textit{ \untertitel }
\end{quote}
selbständig verfasst und keine anderen als die angegebenen Quellen und Hilfsmittel benutzt habe. Die Arbeit wurde bisher keiner anderen Prüfungsbehörde vorgelegt und auch nicht veröffentlicht. \\Die Abgabe der Dokumentation erfolgt gedruckt im Sekretariat der Dualen Hochschule Baden-Württemberg Stuttgart.\\[10ex]

Stuttgart, den \today \\[4ex]


%\rule[-0.2cm]{5cm}{0.5pt} \\	
\begin{figure}[h]
%\includegraphics{Bilder/Unterschrift2.jpg}
\end{figure}
\textsc{\autor} \\[10ex]
 				% Einbinden der eidestattlichen Erklärung
\chapter*{Projektbeschreibung} %*-Variante sorgt dafür, das Abstract nicht im Inhaltsverzeichnis auftaucht


\section{Titel}
   				% Einbinden des Abstracts

\tableofcontents							% Erzeugen des Inhalsverzeichnisses
\printnomenclature[2.0cm]					% Erzeugen des Abkürzungsverzeichnisses
\listoffigures 								% Erzeugen des Abbildungsverzeichnisses 
\listoftables 								% Erzeugen des Tabellenverzeichnisses
\pagebreak


% --------------------------------------------------------------------------------------------
%                    Inhalt der Bachelorarbeit
%---------------------------------------------------------------------------------------------
\pagenumbering{arabic}						% arabische Seitenzahlen für den Hauptteil
\pagestyle{fancy}					
\rmfamily
\noindent

\chapter{Einführung}
\label{cha:Einleitung}


\section{Motivation}
\label{sec:Motivation}

\section{Projektumfeld}
\label{sec:Projektumfeld}

\section{Ziel der Arbeit}
\label{sec:ZielDerArbeit}

\chapter{Text}
\label{cha:text}
\chapter{Reflektion}
\label{cha:reflektion}

\chapter{Ausblick}
\label{cha:ausblick}

%


Für die Vorlage wird das paket \textit{listings} verwendet. \\

\begin{lstlisting} [language=PHP, numbers=left, numberstyle=\tiny, numbersep=10pt]
define('PATH_site', dirname(PATH_thisScript).'/');

if (@is_dir(PATH_site.'typo3/sysext/cms/tslib/')) {
        define('PATH_tslib', PATH_site.'typo3/sysext/cms/tslib/');
} elseif (@is_dir(PATH_site.'tslib/')) {
        define('PATH_tslib', PATH_site.'tslib/');
} else {
      
}
\end{lstlisting}

Quellcodedarstellung 

\begin{verbatim}
user@client:~> texdoc listings
\end{verbatim}

mehr über die Möglichkeiten des Pakets.

------------------------------------------------------\\
\textbf{Notitzen}

20-25 Seiten (Ohne Anhang, inhaltsverzeichniss etc)

Aufgabenstellung

\begin{itemize}
\item[•] Einordnung der Aufgabenstellung in übergeordnete Prozesse/Geschäftsziele
\item[-] Verknüpfung zu Vorlesungsinhalten
\item[-] Praktische Lösung
\item[-] Kritische, inhaltliche Reflexion von Theorie und Praxis
\end{itemize}


------------------------------------------------------\\
\textbf{Form}
Weißes 
Schreibmaschinenpapier  DIN  A  4  (Gewicht  ca.  70  g/qm),  nur  einseitig  beschrieben.  Format:  1,5 
-
zeilig, je Zeile ca. 70 Anschläge, Schriftgröße 12 Punkte
(z.B. bei Arial, b
ei
anderen Schriften ggf. anzupa
s-
sen)
. Randabstand mindestens 2,5 cm allseits.


------------------------------------------------------\\
\textbf{Gliederung der Arbeit}
Die Arbeiten sollten sich in folgende grobe Blöcke unterteilen

\begin{enumerate}
\item Vorspann (Titelblatt, evtl.  Sperrvermerk, Erklärung 
zur Eigenleistung 
[
siehe  2.7
], 
Zusammenfassung
in 
deutsch und englisch, Inhaltsverzeichnis, Abkürzungsverzeichnis, Abbildungs
-
und Tabellen
verzeichnis, 
evtl. Formelgrößen, evtl. Vorwort)

\item Einleitung  (Gegenstand  und  Ziele  der  Arbeit/Aufgabenbeschreibung,  Einführung  in  Thema,  Stand  der 
Technik/Forschung, Motivation der Aufgabenstellung/Vorausblick)
\item 
Hauptteil  (Anforderungsdefinition, 
Anforderungsanalyse,  Lösungsgenerierung,  Lösungsbewertung,  U
m-
setzung) in sinnvollen Gliederungspunkten
\item 
Zusammenfassung und Ausblick
\item 
Literaturverzeichnis
\item 
Anhänge

\end{enumerate}

------------------------------------------------------\\
\textbf{Gestaltung der inhaltlichen Abschnitte/Kapitel}

Einleitung
Die Einleitung soll den Ausgangspunkt der Arbeit umreißen, in kurzer Form zur Problemstellung hinführen 
und das Interesse des Lesers für die Arbeit wecken.
All
gemeine Einleitung ins Thema, keine Unternehmens
-
oder Produktbeschreibungen, Organigramme u.ä., wenn diese nicht direkt zum Thema führen.
Ziele und Vorgehensweise nicht vermischen.
Aufgabenstellung
Die Fragestellung der Aufgabe ist zu präzisieren. 
Insbesondere sind das Umfeld, die vorhandenen Randb
e-
dingungen und Betrachtungsgrenzen darzustellen.
Stand der Technik
Ausgehend von der Aufgabenstellung ist der derzeitige Stand der Technik für die Lösungsfindung zu b
e-
schreiben. Es sind z. B. die Vor
-
und 
Nachteile bisheriger Lösungen bzw. fundamentaler Lösungsprinzipien 
anhand der Literatur darzulegen.
Hauptteil
Der Text soll knapp und klar sein und die wesentlichen Gedanken der Arbeit beinhalten. Ein gewähltes Ve
r-
fahren oder ein bestimmter Lösungsweg muss
begründet werden. Es ist nicht notwendig, alle Vorversuche 
einzeln zu schildern. Bei Versuchen sind Voraussetzungen und Vernachlässigungen sowie die Anordnung, 
Leistungsfähigkeit und Messgenauigkeit der Versuchsanordnung anzugeben.
Die Ergebnisse der Arbe
it sind unter Berücksichtigung der Voraussetzungen ausführlich zu diskutieren und 
mit den bereits bekannten Anschauungen und Erfahrungen zu vergleichen.
Ziel der Arbeit ist es, eindeutige Folgerungen und Richtlinien für die Praxis zu finden.
Zusammenfassun
g 
Aufgabenstellung, Vorgehensweise und wesentliche Ergebnisse werden kurz und präzise dargestellt. 
Die 
Zusammenfassung ist eigenständig verständlich. Länge ca. 1 bis 1,5 Seiten.(Problem, Ziele, Vorgehenswe
i-
se, Ergebnisse und Ausblick).


% ---------------------------- Literaturverzeichnis ----------------------------------------------

\begin{thebibliography}{999999}

%\bibitem[Le01]{levy2001} Autor Name: \emph{Titel des Buches}, New York: Penguin Books, 2001


%\bibitem {Buchtitel} Author: \emph{Buchtitel}, Erscheinungsort: Verlag, Jahr


\end{thebibliography}

% ------------------------------- Anhang ---------------------------------------------------------

\begin{appendix}
\clearpage
\pagenumbering{Roman}						% römische Seitenzahlen für Anhang
\end{appendix}


\end{document}
